\documentclass{article}
\usepackage[UTF8]{ctex}
\usepackage{amsmath,amssymb,amsfonts}
\usepackage{enumitem}
\usepackage{geometry}
\geometry{a4paper, margin=1in}
% === 错题本排版优化（不参与题目标识） ===
\setlength{\parindent}{0pt}
\setlength{\parskip}{0.35em plus 0.10em minus 0.08em}
\setlength{\emergencystretch}{3em}
\setlist{itemsep=0.25em, topsep=0.40em, parsep=0pt}
\allowdisplaybreaks[2]
% === 错题本排版优化结束 ===
\begin{document}
\section*{上海财经大学2025年数学分析真题}
\section*{51．上海财经大学2025年数学分析试卷}

\subsection{题目1}
1．（14 分）计算 \(\lim _{n \rightarrow \infty}(n+1)!\left(e-1-\frac{1}{1!}-\frac{1}{2!}-\cdots-\frac{1}{n!}\right)\) ．

\subsection{题目2}
2．（14 分）计算 \(\lim _{n \rightarrow \infty} n\left(\frac{\sin \frac{\pi}{n}}{n^{2}+1}+\frac{\sin \frac{2 \pi}{n}}{n^{2}+2}+\cdots+\frac{\sin \frac{n \pi}{n}}{n^{2}+n}\right)\) ．

\subsection{题目3}
3．（14 分）求 \(\sum_{n=1}^{\infty} n(n+1) x^{n}\) 的和函数．

\subsection{题目4}
4．（14 分）计算 \(\iint_{D} \sqrt{\frac{1-x^{2}-y^{2}}{1+x^{2}+y^{2}}} \mathrm{~d} x \mathrm{~d} y\) ，其中 \(D=\left\{(x, y) \mid x^{2}+y^{2} \leq 1\right\}\) 在第一象限的部分．

\subsection{题目5}
5．（14 分）计算 \(\int_{L}(x y+y z) \mathrm{d} s\) ，其中 \(L\) 是球面 \(x^{2}+y^{2}+z^{2}=1\) 与平面 \(x+y+z=0\) 的交线．

\subsection{题目6}
6．（10 分）已知 \(\lim _{n \rightarrow \infty}\left(a_{n+2}-a_{n}\right)=a\) ，问 \(\lim _{n \rightarrow \infty} \frac{a_{n}}{n}\) 是否存在，若存在，请给出证明以及它的值．

\subsection{题目7}
7．（10 分）设 \(f(x)\) 在 \([0,3]\) 上可导，\(f(0)+f(1)+f(2)=3, f(3)=1\) ．求证：存在 \(\xi \in(0,3)\) ，使得

\[
f^{\prime}(\xi)=0 .
\]

\subsection{题目8}
8．（10 分）函数列 \(\left\{f_{n}(x)\right\}\) 在 \([a, b]\) 上连续，在 \((a, b)\) 上一致收敛．证明：\(\left\{f_{n}(x)\right\}\) 在点 \(a\) 和 \(b\) 处收敛，并说明极限函数在区间 \([a, b]\) 上连续．

\subsection{题目9}
9．（10 分）已知 \(\sum_{n=1}^{\infty} n a_{n}\) 收敛，判断 \(\sum_{n=1}^{\infty}(n+100) a_{n}\) 的敛散性．

\subsection{题目10}
10．（10 分）设 \(f(x)\) 在 \([0,1]\) 上连续且单调递增，\(g(x)\) 在 \([0,1]\) 上连续且单调递减．证明：

\[
\int_{0}^{1} f(x) g(x) \mathrm{d} x \leq \int_{0}^{1} f(x) \mathrm{d} x \int_{0}^{1} g(x) \mathrm{d} x .
\]

\subsection{题目11}
11．（10 分）设

\[
f(x, y)= \begin{cases}\left(x^{2}+y^{2}\right) \sin \frac{1}{x^{2}+y^{2}}, & (x, y) \neq(0,0) \\ 0, & (x, y)=(0,0)\end{cases}
\]

判断 \(f(x, y)\) 在 \((0,0)\) 处是否可微，偏导数是否连续．

\subsection{题目12}
12．（10 分）设 \(\sum\) 是单位球面 \(x^{2}+y^{2}+z^{2}=1, f\) 是一个连续函数，证明：

\[
\iint_{\sum} f(x+y+z) \mathrm{d} S=2 \pi \int_{-1}^{1} f(\sqrt{3} u) \mathrm{d} u
\]

\subsection{题目13}
13．（10 分）计算 \(f(x, y, z)=\ln x+2 \ln y+4 \ln z\) 在 \(x^{2}+y^{2}+z^{2}=7 R^{2}(x, y, z>0, R>0)\) 下的最大值，并证明对任意的正数 \(a, b, c\) ，都有 \(a b^{2} c^{4} \leq 1024\left(\frac{a+b+c}{7}\right)^{7}\) ．
\end{document}
